\chapter{Formulazione Teorica della Strategia Ottima}
\label{cap:strategia}

Questo capitolo formalizza il problema decisionale dell'operatore di un sistema di accumulo
che opera sul Mercato del Giorno Prima. La formulazione della funzione obiettivo e del
modello fisico \`e impostata; il modello di degrado e la gestione dell'incertezza sono
ancora da sviluppare.

\section{Definizione della Funzione Obiettivo}
\label{sec:obiettivo}

\subsection*{Notazione}

Si consideri una giornata suddivisa in $T$ periodi di mercato indicizzati da
$t = 1, \dots, T$, ciascuno di durata $\Delta_t$ ore ($\Delta_t = 1$ per le aste orarie,
$\Delta_t = 0{,}25$ per quelle a quarto d'ora). Per ogni periodo si definiscono:

\begin{itemize}
  \item $c_t \ge 0$: potenza assorbita in carica nel periodo $t$, in MW;
  \item $s_t \ge 0$: potenza immessa in scarica nel periodo $t$, in MW;
  \item $e_t$: energia immagazzinata alla fine del periodo $t$, in MWh;
  \item $p_t$: prezzo di equilibrio del periodo $t$, in \euromwh{}.
\end{itemize}

\subsection*{Il ricavo e l'effetto di retroazione}

Se l'accumulo fosse di taglia trascurabile rispetto al mercato, il prezzo sarebbe un dato
esogeno e il ricavo giornaliero si scriverebbe semplicemente come
\begin{equation}
\Pi^{\text{price taker}} \;=\; \sum_{t=1}^{T} p_t \,\Delta_t \,(s_t - c_t).
\label{eq:ricavo-pricetaker}
\end{equation}

Questa formulazione, per\`o, ignora il meccanismo che costituisce l'oggetto della tesi: un
accumulo che carica aggiunge domanda e fa \emph{salire} il prezzo, e che scarica aggiunge
offerta e lo fa \emph{scendere}. Il prezzo non \`e quindi un dato ma una funzione della
decisione stessa:
\begin{equation}
p_t \;=\; \pi_t\!\left(c_t, s_t\right),
\label{eq:prezzo-endogeno}
\end{equation}
dove $\pi_t$ \`e la funzione che restituisce il prezzo di equilibrio del periodo $t$ dopo
aver inserito nelle curve d'asta la domanda addizionale $c_t$ e l'offerta addizionale $s_t$.
Questa funzione non ha forma analitica chiusa: \`e definita dall'algoritmo di
clearing~\eqref{eq:clearing} applicato alle curve ricostruite dai dati, ed \`e implementata
nel modulo descritto nella Sezione~\ref{sec:setup}.

Il problema dell'operatore diventa quindi
\begin{equation}
\max_{\{c_t,\, s_t\}} \quad
\Pi \;=\; \sum_{t=1}^{T} \pi_t\!\left(c_t, s_t\right) \Delta_t \left(s_t - c_t\right)
\;-\; \Phi\!\left(\{c_t, s_t\}\right),
\label{eq:obiettivo}
\end{equation}
dove $\Phi$ rappresenta il costo di degrado associato al profilo di utilizzo
(Sezione~\ref{sec:degrado}).

\paragraph{Osservazione.} La retroazione agisce \emph{contro} l'operatore in entrambe le
direzioni: comprando fa salire il prezzo di acquisto, vendendo fa scendere quello di vendita.
Il margine di arbitraggio si riduce quindi al crescere della taglia dell'accumulo, ed \`e
questo il meccanismo che rende il dimensionamento un problema non banale — al di sopra di una
certa capacit\`a, l'accumulo distrugge da s\'e la propria fonte di reddito. La misura di
questo effetto \`e riportata nella Sezione~\ref{sec:esecuzione}.

\section{Modello fisico della batteria}
\label{sec:modello-fisico}

Il bilancio energetico dell'accumulo lega lo stato di carica fra periodi consecutivi:
\begin{equation}
e_t \;=\; e_{t-1} \;+\; \eta^{c}\, c_t\, \Delta_t \;-\; \frac{s_t\, \Delta_t}{\eta^{s}},
\qquad t = 1, \dots, T,
\label{eq:bilancio}
\end{equation}
dove $\eta^{c}$ ed $\eta^{s}$ sono i rendimenti di carica e di scarica; il loro prodotto
$\eta = \eta^{c}\eta^{s}$ \`e il rendimento di ciclo completo.

Il problema \`e soggetto ai vincoli
\begin{align}
0 \;\le\; e_t \;&\le\; E^{\max}, \label{eq:vincolo-energia}\\
0 \;\le\; c_t \;&\le\; P^{\max}, \label{eq:vincolo-carica}\\
0 \;\le\; s_t \;&\le\; P^{\max}, \label{eq:vincolo-scarica}\\
c_t \cdot s_t \;&=\; 0, \label{eq:vincolo-esclusione}
\end{align}
dove $E^{\max}$ \`e la capacit\`a energetica in MWh e $P^{\max}$ la potenza nominale in MW.
Il loro rapporto $E^{\max}/P^{\max}$ \`e la \emph{durata} dell'accumulo, espressa in ore, ed
\`e uno dei parametri di dimensionamento che la tesi si propone di ottimizzare.

Il vincolo~\eqref{eq:vincolo-esclusione} impedisce di caricare e scaricare simultaneamente.
\`E un vincolo di complementarit\`a, non lineare; nella pratica pu\`o essere omesso quando i
prezzi sono positivi, perch\'e caricare e scaricare nello stesso periodo sarebbe comunque
non conveniente, ma va reintrodotto in presenza di prezzi negativi, che sui dati esaminati
sono ammessi fino a $-500$~\euromwh{}.

% DA COMPLETARE: scelta dei valori dei parametri fisici.
% Da fissare, con riferimento a fonti tecniche: rendimento di ciclo completo (tipicamente
% 0,85-0,92 per batterie agli ioni di litio), profondita' di scarica ammessa, eventuali
% vincoli di rampa, autoscarica. Va inoltre deciso se imporre e_T = e_0 (ciclo giornaliero
% chiuso), scelta che semplifica il confronto fra giornate.

\section{Modellazione del degrado}
\label{sec:degrado}

% DA COMPLETARE: modello di degrado della batteria.
% Elementi da trattare:
%   - distinzione fra degrado calendariale (funzione del tempo e dello stato di carica
%     medio) e degrado ciclico (funzione del numero e della profondita' dei cicli);
%   - scelta di una formulazione trattabile: costo marginale costante per MWh ciclato,
%     oppure conteggio dei cicli equivalenti con metodo rainflow;
%   - traduzione del degrado in un costo monetario Phi da inserire nella funzione
%     obiettivo (eq. 3.4), tipicamente come costo di sostituzione ripartito sui cicli
%     di vita utile;
%   - implicazione economica: il degrado introduce un costo opportunita' che rende non
%     conveniente ogni arbitraggio il cui differenziale di prezzo non lo copra, e quindi
%     riduce il numero di cicli ottimali rispetto al caso senza degrado.

\section{Gestione dell'incertezza}
\label{sec:incertezza}

% DA COMPLETARE: trattamento dell'incertezza sui prezzi.
% Il punto di partenza e' che la strategia ottimizzata sui prezzi realizzati costituisce
% un limite superiore irraggiungibile (perfect foresight): serve come termine di confronto,
% non come strategia realistica.
% Elementi da trattare:
%   - quantificazione del divario fra strategia con previsione perfetta e strategie
%     realizzabili sulla base dell'informazione disponibile alla chiusura del MGP;
%   - possibili impostazioni: ottimizzazione stocastica su scenari di prezzo, oppure
%     regole di decisione basate su previsioni puntuali con analisi dell'errore;
%   - collegamento con la Sezione 2.4 sull'aleatorieta' del prezzo.
